% =====================================================================
% 1. DEVELOPING MACHINE LEARNING MODELS
% =====================================================================
\section{Developing Machine Learning models}

\subsection{Problem description}

Diffuse gliomas are primary tumours of the central nervous system, and their standard-of-care
treatment ``is predicated on tumor grade'' \cite{tasci2022}. For decades they were graded from
histology alone, but the histological classification of
lower-grade gliomas ``suffers from high intraobserver and interobserver variability and does not
adequately predict clinical outcomes'' \cite{tcga2015lgg}. The 2016 and 2021 World Health
Organization (WHO) classifications therefore integrate molecular markers with histology
\cite{louis2016,louis2021}: in the 2021 scheme, the common adult-type diffuse gliomas comprise ``only 3 types'':
astrocytoma, IDH-mutant (CNS WHO grades 2--4), oligodendroglioma, IDH-mutant and
1p/19q-codeleted (grades 2--3), and glioblastoma, IDH-wildtype (grade 4) \cite{louis2021}.
The grade matters because prognosis differs widely: in the TCGA lower-grade glioma cohort,
patients with an IDH mutation and 1p/19q codeletion had a median survival of 8.0 years, and those
with an IDH mutation without codeletion of 6.3 years \cite{tcga2015lgg}, whereas glioblastoma,
which accounts for 14.3\% of all primary brain and other CNS tumours in the United States and is
diagnosed at a median age of 65 years, has a five-year relative survival of only 6.8\%
\cite{ostrom2021}.

The molecular markers relevant to this project are point mutations. Gliomas without IDH1 or
IDH2 mutations tend to carry other alterations: among anaplastic astrocytomas and glioblastomas
with wild-type IDH genes, \cite{yan2009} found few TP53 mutations (18\%) and frequent alterations
of PTEN, EGFR, CDKN2A or CDKN2B (74\%). In routine practice these alterations are detected with
immunohistochemistry, sequencing and copy-number assays; for example, fluorescence \emph{in situ}
hybridisation detects gains of chromosome 7 and losses of chromosome 10, the 1p/19q codeletion,
homozygous deletions such as CDKN2A/B, and gene amplifications \cite{sahm2023}.

This motivates the questions of this project, stated as hypotheses before any model was trained:
\begin{enumerate}[label=\textbf{H\arabic*.}, leftmargin=*]
  \item The mutation status of 20 genes, together with age, sex and race, contains enough
        information to separate lower-grade glioma (LGG) from glioblastoma (GBM) well above a trivial
        classifier, and simple interpretable models perform close to complex ones.
  \item Without using the grade, the same profile explains part of the variability of the age at
        diagnosis, mostly through the markers that also define the grade.
  \item Without any label, the mutation profiles form groups that are related to, but not
        identical with, the grade.
\end{enumerate}
After the analysis of H1--H3, two complementary questions were added, and they are reported as such
rather than as hypotheses fixed in advance:
\begin{enumerate}[label=\textbf{Q\arabic*.}, leftmargin=*, start=4]
  \item Can the same profile go beyond the grade and recognise the histological family
        (astrocytoma, oligodendroglioma, mixed glioma, glioblastoma)?
  \item Under explicit causal assumptions, how much of the association between an IDH mutation and a
        younger age at diagnosis goes through the grade?
\end{enumerate}
The objective is not to beat published accuracies but to apply and critically compare the
models of the course under a validation protocol that gives honest estimates.

\subsection{Machine Learning approach}

The three hypotheses map onto the three learning tasks required in the course:
\begin{itemize}
  \item \textbf{Classification} (H1): target \texttt{Grade} (LGG = 0, GBM = 1), predictors age,
        sex, race and 20 mutations. Models: penalised logistic regression, naive Bayes, a pruned
        classification tree, random forest, gradient boosting (XGBoost) and support vector machines.
  \item \textbf{Regression} (H2): target \texttt{Age\_\allowbreak{}at\_\allowbreak{}diagnosis} (years), predictors sex, race
        and 20 mutations, deliberately \emph{without} the grade. Models: multiple linear regression,
        Ridge, Lasso, a regression tree, random forest, XGBoost and support vector regression.
  \item \textbf{Multiclass classification} (Q4): target the histological family derived from the
        diagnosis (four classes), same predictors as H1 and \emph{without} the grade. Models:
        multinomial logistic regression, a pruned tree, random forest and XGBoost. The winning models of
        both classification tasks are then explained with SHAP values.
  \item \textbf{Causal analysis} (Q5): a causal diagram and a mediation decomposition with linear
        regression, with the grade as mediator between IDH and the age at diagnosis.
  \item \textbf{Unsupervised learning} (H3): PCA, multidimensional scaling and t-SNE with a
        Jaccard distance, and $k$-means and hierarchical clustering of the 20 mutation indicators;
        the grade is used only afterwards, to describe the clusters.
\end{itemize}
Every supervised model is evaluated with the same nested cross-validation and compared with a
trivial baseline (Section~\ref{sec:prdl-setup}); the tools and the implementation are described
in the companion MLLB report.
